% Compile with XeLaTeX or LuaLaTeX and an installed ko.TeX distribution.
\documentclass[11pt]{article}
\usepackage{kotex}
\usepackage{amsmath,amssymb,booktabs}
\usepackage[margin=1in]{geometry}
\usepackage[hidelinks]{hyperref}
\newcommand{\KL}{\operatorname{KL}}
\newcommand{\NLL}{\operatorname{NLL}}
\newcommand{\norm}[1]{\left\lVert #1\right\rVert}
\newcommand{\mean}{\operatorname{mean}}
\newcommand{\pos}[1]{\left[#1\right]_+}
\title{JLZ v12: 공유 native 예산 아래 층별 편집 효율 배분}
\author{한국어 방법 명세}
\date{2026년 10월 4일}
\begin{document}
\maketitle

\paragraph{범위와 주장.}
JLZ v12는 모든 쓰기 층의 subject 변화를 한 forward 경로에서 공동으로
계획하고, 층별 상대 변화의 norm 합에 하나의 예산을 부여한다.
계획의 목적은 native 편집 손실을 개선하는 것이며, 층별 보존 비용을
명시적으로 최적화하는 것은 아니다. 이후 native ridge writer가 층별
virtual 목표와 현재 hidden 상태의 차이를 순차적으로 쓴다.
계획 예산은 virtual 주입량에 대한 제약이다. 실제 가중치 변화 또는
모든 토큰의 hidden 변화에 같은 상한을 보장하지 않는다.
유한 횟수의 EfficiencyAdam은 배분을 구하는 휴리스틱이며, 아래 KKT
조건은 그 결과의 진단 기준이다. 수렴이나 KKT 만족을 주장하지 않는다.

\section{고정 entry 상태와 층별 변수}
순차 batch $t$의 entry 모델을 $W_t$, 기존 key-history covariance를
$H_{t,l}$이라 한다. 실제 요청 수 $B$, 쓰기 층 집합 $\mathcal L$,
층별 block 출력 차원 $d_l$, writer 입력 차원 $k_l$은 실행 profile이
정한다. 마지막 partial batch에도 실제 $B$를 사용한다.
층별 $d_l$과 $k_l$이 같다고 가정하지 않는다.
native z 층 $l_\star$와 loss readout 층도 profile이 정한다.
Anchor 층이 $\mathcal L$ 밖에 있으면 그 entry 출력은 별도로 측정한다.

요청 $r$의 canonical entry subject에서 full transformer-block 출력을
$h^0_{lr}$라 하고 다음 anchor를 고정한다.
\begin{equation}
 a_{lr}=\norm{h^0_{lr}}_2>0,\qquad
 a^\star_r=\norm{h^0_{l_\star r}}_2>0.
 \label{eq:anchors}
\end{equation}
$W_t$, $H_{t,l}$, anchor, entry KL teacher $p^0_{K,r}$는 계획 동안
변하지 않는다. Anchor가 0이거나 유한하지 않으면 기술적 실패로
기록하며, 숨겨진 floor로 목적을 바꾸지 않는다.
Anchor norm은 native activation 및 norm reduction dtype으로 entry에서
한 번 계산해 고정한다. 현재 profile은 FP32이며, 사영의 FP64 계산과
혼동하여 anchor norm을 임의로 FP64에서 다시 계산하지 않는다.

각 $l\in\mathcal L$에서 최적화하는 상대 변화와 실제 주입량은
\begin{equation}
 u_{lr}\in\mathbb R^{d_l},\qquad
 \delta_{lr}=a_{lr}u_{lr},\qquad u^{(0)}_{lr}=0.
 \label{eq:parameterization}
\end{equation}
요청의 모든 native rewrite 및 KL context에서 subject lookup 위치의
block 출력에 동일한 $\delta_{lr}$를 더한다. 모든 층의 주입을 같은
forward에 적용하므로 아래층 변화가 위층 hidden과 최종 손실에 미치는
경로를 미분한다. 층마다 독립적인 task loss를 두지 않는다.
모델 가중치는 optimizer 변수가 아니지만 주입량까지의 gradient 경로는
유지한다. Subject 위치, token 및 target 위치, padding과 position 의미,
rewrite/KL context와 readout은 native 정의를 따른다.

계획의 주입 위치는 full block 출력이고 writer key는 선택한
down-projection의 입력이다. 모델 adapter는 두 위치의 국소 가산 관계와
weight orientation을 확인해야 한다. 지원하지 않는 구조에 조용히
대체 규칙을 적용하지 않는다. 평가용 paraphrase나 neighborhood 문장을
계획, stopping 또는 계수 선택에 사용하지 않는다.

\section{편집 목적과 공유 예산}
요청마다 다음 constrained objective를 사용한다.
\begin{align}
 F_r(u_r)
 &=\NLL_r^v+\lambda_K\KL(p^v_{K,r}\Vert p^0_{K,r}),
 \label{eq:task}\\
 \beta_r&=\frac{\lambda_n}{a^\star_r},\qquad
 b_r(u_r)=\sum_{l\in\mathcal L}\norm{u_{lr}}_2,\\
 J_r(u_r)&=F_r(u_r)+\beta_r b_r(u_r),
 & b_r(u_r)&\le c.
 \label{eq:objective}
\end{align}
NLL은 native target-token 및 context reduction을 유지하고, KL은
full vocabulary에서 \emph{current-to-entry} 방향으로 계산한다.
Native key 추출의 context 평균과 NLL의 context 평균은 서로 다른
reduction일 수 있으므로 임의로 통일하지 않는다.
Decay는 norm의 제곱이 아닌 norm의 합이다.
층마다 상대 변화 한 단위에 동일한 가격 $\beta_r$를 적용하며, 별도
allocation coefficient나 용량ㆍold-costㆍreference loss를 추가하지 않는다.

Group-$\ell_1$ 예산 $b_r\le c$는 다층 계획에 native 크기의 상대 변화
예산 하나를 부여한다. 각 층에 $c$씩 허용하는 독립 clamp와 다르다.
이 제약은 일부 층에 0을 배분하는 해를 허용하지만, 특정 층의 선택이나
앞층 집중을 강제하지 않는다. 계층별 hidden 공간은 서로 다르므로
$b_r$를 실제 모델 출력 변화의 norm으로 해석하지 않는다.
사영으로 0이 된 층도 다음 update에서 다시 예산을 받을 수 있다.
Moment 삭제, 영구 pruning 또는 비제로 층만의 재정규화를 하지 않는다.
$b_r=c$인 경계에서 decay는 $\beta_r c$로 일정하므로, 이 항 자체가
경계 위 층간 선호를 정하지 않는다. 유한 update의 배분에는 optimizer의
방향과 이력도 영향을 미친다.

요청끼리 결합하는 목적 항이나 예산은 없다. 같은 batch로 계산하더라도
$J_r$, anchor, moment, 예산, stopping과 candidate counter는 요청별로
독립이다. Batch 평균을 보고용으로 계산할 수 있으나, optimizer에 넣는
gradient는 batch 크기에 따라 축소되지 않은 $\nabla_{u_r}J_r$이다.

\section{한계 효율의 KKT 해석과 한계}
$F_r$가 미분 가능한 점에서 norm 제약의 multiplier를 $\mu_r\ge0$라
하고, $u_{lr}\ne0$이면 $e_{lr}=u_{lr}/\norm{u_{lr}}_2$라 하자.
정지점에서 필요한 KKT 조건은
\begin{align}
 \nabla_{u_{lr}}F_r+(\beta_r+\mu_r)e_{lr}&=0,
 &&u_{lr}\ne0,\label{eq:kkt-active}\\
 \norm{\nabla_{u_{lr}}F_r}_2&\le\beta_r+\mu_r,
 &&u_{lr}=0,\label{eq:kkt-inactive}\\
 b_r&\le c,\qquad \mu_r\ge0,\qquad
 \mu_r(b_r-c)=0.
 \label{eq:kkt-feasibility}
\end{align}
활성 층에서는 gradient의 방향도 $-e_{lr}$와 일치해야 하므로 norm만
같은 것으로 stationarity가 성립하지 않는다. 이 조건에서 task-gradient
norm은 공통 가격 $\beta_r+\mu_r$와 같고, 비활성 층은 이를 넘지 않는다.
예산이 남으면 $\mu_r=0$일 수 있으며 공통 가격은 여전히 $\beta_r$이다.

이 해석은 native 편집 목적 $F_r$의 상대 변화 단위당 효율에 관한
것이다. 실제 writer의 보존 비용 또는 locality의 한계 효율과 동일하지
않다. 신경망 목적은 일반적으로 비볼록이며, KKT는 global optimum이나
유한 optimizer 궤적의 성질을 보장하지 않는다.
실험에서는 활성 층의 식~\eqref{eq:kkt-active} 잔차, 비활성 층의
식~\eqref{eq:kkt-inactive} 위반, feasibility와 complementarity를
함께 기록한다. Multiplier 추정법과 diagnostic active tolerance도
기록하되, 이 진단을 stopping, admission, candidate 선택 또는 추가
update의 조건으로 사용하지 않는다.
진단용으로 예산 내부에서는 $\widehat\mu_r=0$, 경계에서 활성 층
집합 $\mathcal A_r$가 비어 있지 않으면
\begin{equation}
 \widehat\mu_r=\max\left\{0,\frac1{|\mathcal A_r|}
 \sum_{l\in\mathcal A_r}\bigl[-\langle\nabla_{u_{lr}}F_r,e_{lr}\rangle
 -\beta_r\bigr]\right\}
 \label{eq:mu-diagnostic}
\end{equation}
를 사용한다. Active 및 경계 tolerance는 telemetry schema에 고정하고
raw norm도 저장한다. 진단상의 zero 분류가 실제 변수를 바꾸지는 않는다.
Native의 stop-before-backward 순서를 유지하므로 종료 후보에는 추가
backward를 하지 않는다. 그 후보의 KKT는
\texttt{NOT\_EVALUATED(NO\_BACKWARD\_TERMINAL)}로 기록하고,
gradient를 계산한 정확한 후보에 대해서만 진단한다. 이전 후보의
잔차를 terminal 후보의 잔차로 대신하지 않는다.
주 profile은 $c>0$이다. CPU 검증용 퇴화 경우 $c=0$에서는 $u=0$만
가능하며, $\mu_r=\max\{0,\max_l\norm{\nabla_{u_{lr}}F_r}_2-\beta_r\}$를
사용하면 비활성 조건을 만족하는 multiplier를 구성할 수 있다.

\section{EfficiencyAdam: 유한 예산의 배분 휴리스틱}
평가한 후보에서 아직 종료하지 않은 요청의 gradient를
\begin{equation}
 g^{(q)}_{lr}=g^{(q),F}_{lr}+\beta_r e_{lr},\qquad
 g^{(q),F}_{lr}=\nabla_{u_{lr}}F_r,\qquad
 e_{lr}=\begin{cases}
 u_{lr}/\norm{u_{lr}}_2,&u_{lr}\ne0,\\
 0,&u_{lr}=0
 \end{cases}
 \label{eq:gradient}
\end{equation}
로 계산한다. 모델 backward는 $F_r$에 대해서만 수행하고 norm 항의
gradient는 해석적으로 더한다. 이는 선택한 zero-norm subgradient를
사용한 $\nabla J_r$와 같으며, task gradient 진단을 위해 별도 model
backward를 추가하지 않는다. Norm 안에 smoothing 상수를 추가하지 않는다.
요청별 update 번호 $q=1,2,\ldots$에 대해 좌표별 Adam moment와
층별 squared-gradient 평균을 다음과 같이 갱신한다.
\begin{align}
 m^{(q)}_{lr}&=\rho_1m^{(q-1)}_{lr}+(1-\rho_1)g^{(q)}_{lr},\\
 v^{(q)}_{lr}&=\rho_2v^{(q-1)}_{lr}+(1-\rho_2)(g^{(q)}_{lr})^{\odot2},\\
 s^{(q)}_{lr}&=\rho_2s^{(q-1)}_{lr}
 +(1-\rho_2)\frac1{d_l}\sum_{i=1}^{d_l}(g^{(q)}_{lr,i})^2,\\
 \widehat m^{(q)}_{lr}&=\frac{m^{(q)}_{lr}}{1-\rho_1^q},\qquad
 \widehat v^{(q)}_{lr}=\frac{v^{(q)}_{lr}}{1-\rho_2^q},\qquad
 \widehat s^{(q)}_{lr}=\frac{s^{(q)}_{lr}}{1-\rho_2^q},\\
 \overline s^{(q)}_r&=\frac1{|\mathcal L|}
                      \sum_{l\in\mathcal L}\widehat s^{(q)}_{lr},\qquad
 \gamma^{(q)}_{lr}=\sqrt{\frac{\widehat s^{(q)}_{lr}}
                                   {\overline s^{(q)}_r}}.
 \label{eq:gamma}
\end{align}
Moment는 0에서 시작한다. $\overline s_r=0$이면 모든 $\gamma_{lr}=1$로
정의한다. 평균은 각 층의 좌표 수로 먼저 정규화하고
층 사이에서는 산술평균을 사용한다. 따라서 층별 폭이 달라도 식이
정의되지만 RMS 비를 gradient norm 비와 동일시할 수는 없다.

Native learning rate와 epsilon을 $\eta_{\rm native}$,
$\epsilon_{\rm native}$라 할 때
\begin{align}
 \eta_r&=\frac{\eta_{\rm native}}{a^\star_r},&
 \epsilon_{u,r}&=a^\star_r\epsilon_{\rm native},\\
 \widetilde u_{lr}&=u_{lr}
 -\eta_r\gamma_{lr}
 \frac{\widehat m_{lr}}{\sqrt{\widehat v_{lr}}+\epsilon_{u,r}}.
 \label{eq:adam}
\end{align}
나눗셈과 제곱근은 좌표별 연산이다. 기본 profile은
$\eta_{\rm native}=0.1$, $(\rho_1,\rho_2)=(0.9,0.999)$,
$\epsilon_{\rm native}=10^{-8}$을 사용한다.
Explicit norm decay 외 optimizer weight decay, warm-up, momentum 재시작을
추가하지 않는다. Plan과 moment는 현재 profile에서 FP32로 저장한다.

$\gamma_{lr}$는 좌표별 Adam 정규화가 줄이는 층 사이 gradient 크기
차이를 반영하기 위한 휴리스틱이다. KKT multiplier를 계산하거나
constrained objective를 정확히 푸는 단계가 아니다. 첫 update에서도
좌표별 epsilon, gradient 방향, 차원과 뒤이은 사영 때문에 최종 배분
비율이 gradient norm 비와 정확히 같다고 주장하지 않는다.

\section{정확한 Euclidean 사영과 요청별 종료}
매 update 뒤 다음 사영을 수행한다.
\begin{equation}
 u^{\rm next}_r=\mathop{\arg\min}_{w}
 \frac12\sum_{l\in\mathcal L}\norm{w_l-\widetilde u_{lr}}_2^2
 \quad\text{subject to}\quad
 \sum_{l\in\mathcal L}\norm{w_l}_2\le c.
 \label{eq:projection}
\end{equation}
$n_l=\norm{\widetilde u_{lr}}_2$라 하면, $\sum_l n_l\le c$일 때
후보를 그대로 둔다. 그 외에는
\begin{equation}
 \sum_l\pos{n_l-\tau_{{\rm proj},r}}=c,\qquad
 u^{\rm next}_{lr}=\begin{cases}
 \displaystyle\frac{\pos{n_l-\tau_{{\rm proj},r}}}{n_l}\widetilde u_{lr},
     & n_l>0,\\
 0,& n_l=0
 \end{cases}
 \label{eq:threshold}
\end{equation}
인 scalar threshold를 구한다. 이는 사영 subproblem의 정확한 해이다.
Adam moment $m,v,s$는 사영 뒤에도 유지한다.
이 사영은 Euclidean 거리 기준이며 Adam preconditioner의 거리 기준
사영이 아니다. $\tau_{{\rm proj},r}$는 한 update의 사영 임계값이고,
식~\eqref{eq:kkt-active}의 최종 목적 multiplier $\mu_r$ 또는 공통
가격 $\beta_r+\mu_r$와 같다고 둘 수 없다.
사영 계산은 FP64로 수행하고 결과를 FP32 plan으로 cast한 뒤 실제
저장값의 feasibility를 검사한다. 고정 tolerance를 넘으면 기술적 실패로
기록하며, 숨겨진 반복 축소나 추가 배분 정책으로 보정하지 않는다.

요청마다 최대 25회 후보 평가와 24회 optimizer update를 허용한다.
첫 평가는 $u=0$이며 다음 순서를 따른다.
\begin{enumerate}
 \item 현재 후보의 $F_r$, decay와 전체 $J_r$를 평가하고 유한성을 확인한다.
 \item $J_r<\varepsilon_{\rm stop}$이면 그 요청을 종료한다. 첫 평가에서 만족하면
       $\delta_{lr}=0$인 zero-step 결과를 채택한다.
 \item 이미 25회 평가했다면 현재 후보를 채택한다.
 \item 그 외에는 gradient, EfficiencyAdam update와 사영을 한 번 수행한
       다음 후보를 평가한다.
\end{enumerate}
종료한 요청의 후보, moment와 counter는 고정하며 다른 요청만 계속한다.
현재 profile의 $\varepsilon_{\rm stop}$은 $0.05$이다.
Batch 평균이나 NLL만으로 종료하지 않는다. 정상 종료 시 마지막으로
평가한 유한 후보를 사용하고, 평가하지 않은 post-update 후보나
평가 지표가 좋은 과거 후보를 선택하지 않는다. 비유한 objective,
gradient 또는 상태는 기술적 실패이며 조용히 과거 후보로 복구하지
않는다. KKT 잔차가 크다는 이유로 추가 step을 허용하지 않는다.

\section{최종 virtual 목표와 causal ridge write}
요청별로 채택한 후보를 고정하고, entry 모델의 canonical subject에서
모든 선택된 주입을 활성화한 최종 목표를 사용한다. 모든 loss 평가
forward에서 canonical post-injection 출력을 함께 capture하므로,
종료한 후보의 목표를 얻기 위한 추가 forward는 없다.
\begin{equation}
 z^v_{lr}=h^{v,-}_{lr}+\delta_{lr}.
 \label{eq:virtual-target}
\end{equation}
$h^{v,-}_{lr}$는 아래층 virtual 주입을 이미 거친 층 $l$의 block 출력이며
자기 층 주입 직전의 값이다. 따라서 $z^v_{lr}$는 자기 층 주입까지
포함한 \emph{post-injection} full block 출력이다.
종료 시 해당 평가에서 capture한 값을 고정하며, 이전 후보의 목표를
재사용하거나 종료 뒤 새 후보를 만들지 않는다.

$(W_t,H_t)$에 기반한 write transaction 안에서 모든
$l\in\mathcal L$을 native causal 순서, 즉 아래층부터 처리한다.
\begin{enumerate}
 \item 이미 앞층 write를 반영한 현재 모델에서 native
       \texttt{compute\_ks}와 context reduction으로 실제 key
       $K_l\in\mathbb R^{k_l\times B}$를 새로 추출한다.
 \item 같은 현재 모델에서 canonical subject의 full block 출력
       $h^{\rm cur}_{lr}$를 다시 측정한다.
 \item $R_l=[r_{l1},\ldots,r_{lB}]$,
       $r_{lr}=z^v_{lr}-h^{\rm cur}_{lr}$로 residual을 구성한다.
 \item Native ridge solve와 update를 적용한다.
       \begin{align}
       P_l&=(\lambda_C C_{0,l}+H_{t,l}+K_lK_l^\top)^{-1}K_l,\\
       U_l&=R_lP_l^\top.
       \label{eq:writer}
       \end{align}
       역행렬 표기는 선형 solve를 뜻한다. Native weight orientation,
       dtype cast와 addition을 적용한 다음 다음 층으로 진행한다.
\end{enumerate}
남은 층 수 divisor, inverse-realization 보정 또는 별도의 pulse를
사용하지 않는다. Entry key를 actual writer key 대신 재사용하지 않는다.
계획 $\delta_{lr}=0$인 요청ㆍ층도 write 대상이며, 한 층의 모든 계획
변화가 0이어도 inherited mismatch가 있으면 그 층의 write는 0이 아닐
수 있다. 계획 active set은 writer eligibility gate가 아니다.

모든 층의 write를 마친 뒤 final 모델에서 native mean history key
$K_{{\rm hist},l}$를 다시 추출하고, 각 층에 한 번씩
\begin{equation}
 H_{t+1,l}=H_{t,l}+K_{{\rm hist},l}K_{{\rm hist},l}^{\top}
 \label{eq:history}
\end{equation}
를 native history dtype 및 저장 규칙으로 append한다.
중간층 commit 전에 history를 갱신하지 않는다. Key 재사용은 final-key
parity를 사전에 확인한 경우에만 허용한다. 실패 시 transaction을
되돌리고 비선택 weight가 유지되는지 확인한다.
Zero-step 요청도 batch writer의 원래 key 및 target 열에 남는다.
성공한 요청만 선택해 append하거나 중복 요청을 제거하지 않는다.

\section{계획, inherited mismatch와 실제 local 작용의 구별}
각 writer residual은 정확히
\begin{equation}
 r_{lr}=\underbrace{\delta_{lr}}_{\text{자기 층 계획}}
 +\underbrace{\bigl(h^{v,-}_{lr}-h^{\rm cur}_{lr}\bigr)}_
                  {\text{inherited mismatch}}
 \label{eq:residual-decomposition}
\end{equation}
로 나뉜다. 두 번째 항은 아래층 virtual 주입과 실제 weight write가
만든 상태 차이이다. 이를 항상 양의 크기를 가진 ``부족분''으로
해석할 수는 없다. 방향과 norm을 별도로 기록해야 한다.
Ridge가 요청끼리 결합하므로 한 요청의 계획이나 residual 열이 0이어도
다른 요청에 대한 write가 그 요청의 key에 작용할 수 있다.

계획 share와 canonical subject에서 실제 local write의 share를
구분한다. $\Delta W_l^{\rm applied}$는 native cast와 실제 addition
후 측정한 weight 차이를 canonical 출력 orientation으로 나타낸 것이고,
$k^{\rm can}_{lr}$는 층 $l$ write 직전 현재 모델의 canonical writer
입력이다. Adapter의 국소 가산 관계가 성립할 때
\begin{align}
 y^{\rm local}_{lr}&=\Delta W_l^{\rm applied}k^{\rm can}_{lr},\\
 q^{\rm plan}_{lr}&=\frac{\norm{\delta_{lr}}_2}{a_{lr}},&
 q^{\rm local}_{lr}&=\frac{\norm{y^{\rm local}_{lr}}_2}{a_{lr}},\\
 \pi^{\rm plan}_{lr}&=\frac{q^{\rm plan}_{lr}}
                           {\sum_j q^{\rm plan}_{jr}},&
 \pi^{\rm local}_{lr}&=\frac{q^{\rm local}_{lr}}
                            {\sum_j q^{\rm local}_{jr}}.
 \label{eq:shares}
\end{align}
계획 분모가 0이면 0벡터와 \texttt{NO\_PLANNED\_EDIT}를 기록하며
이를 통상적인 합 1의 share로 해석하지 않는다. 실제 local 작용의
분모가 0이면 \texttt{null}과 \texttt{ZERO\_REALIZED\_TOTAL}을 기록한다.
Native mean key에서의 solve상 작용 $U_lK_l$과 canonical local 작용은
key와 cast가 달라 서로 구분한다. 누적 block-output 변화에는 아래층
효과가 섞이므로 이를 local realized share로 대신하지 않는다.
실제 local 작용에는 식~\eqref{eq:residual-decomposition}의 두 항이
함께 기여한다. 따라서 이를 순수한 자기 층 계획의 실현율로도
해석하지 않는다.

계획 주입은 subject 위치에 한정되지만 실제 weight write는 다른 토큰과
문장에도 작용하고, ridge는 계획을 정확히 실현하지 않을 수 있다.
공유 예산만으로 write norm, 실제 상대 local 크기의 합
$\sum_l q^{\rm local}_{lr}$, locality 또는
fit-to-write task parity를 보장하지 않는다. Locality의 이득과 lifelong
retention은 별도 실험으로 판단한다.

\section{제한적인 단일 층 native 환원}
$\mathcal L=\{l_\star\}$이고 같은 canonical anchor, 입력, readout,
reduction, dtype, native optimizer와 stopping 설정을 사용할 때에만
다음 planning 대응을 주장한다.
\begin{align}
 u_r&=\delta_r/a^\star_r,\\
 \beta_r\norm{u_r}_2
 &=\lambda_n\frac{\norm{\delta_r}_2}{(a^\star_r)^2},&
 \norm{u_r}_2\le c
 &\ \Longleftrightarrow\ \norm{\delta_r}_2\le c a^\star_r.
\end{align}
비영 gradient가 있는 단일 층에서는 $\gamma=1$이다.
$g_u=a^\star g_\delta$, $\widehat m_u=a^\star\widehat m_\delta$,
$\widehat v_u=(a^\star)^2\widehat v_\delta$이므로
식~\eqref{eq:adam}의 $\eta_r=\eta_{\rm native}/a^\star_r$와
$\epsilon_{u,r}=a^\star_r\epsilon_{\rm native}$를 함께 사용하면
동일한 native Adam update를 $\delta$ 단위로 얻는다. 모두 zero-gradient인
경우도 update가 0이라는 의미에서 대응한다. 유한 정밀도에서는 명시한
tolerance로 loss, optimizer step과 clamp의 수치 parity를 확인해야 한다.

이는 native z 층에서의 단일 변수 planning에 대한 제한적인 환원이다.
다층 MEMIT-H writer 전체 또는 다른 native edit pipeline 전체의
동등성을 뜻하지 않는다. $\mathcal L=\{l\}$에서 $l\ne l_\star$이면
공통 anchor $a^\star_r$ 때문에 해당 층의 native decay와 learning rate가
자동으로 일치하지 않는다. 첫 층에 계획을 고정한 비교군은 배분의
효과를 검사하는 ablation이며, FE의 전체 절차와 정확히 같다고 부르지
않는다.

\section{Profile, 검증과 기록}
방법의 수식은 모델, layer 번호와 benchmark에 독립적이다. 하나의 예시
profile은 쓰기 층 L4--L8, native z anchor L8, loss readout 31,
$\lambda_K=0.0625$, $\lambda_n=0.5$, $c=0.75$,
$\lambda_C=15000$을 사용한다. 이 수치는 방법의 정의에 hard-code할
상수가 아니며 다른 profile은 native 설정과 adapter를 명시해야 한다.
이 profile의 geometry는 FP64, history는 native CPU FP32를 사용한다.
기존 2,000 edits와 BS100$\times$20 범위는 후속 실험 workload로
유지하되, 본 명세가 해당 실행이나 자원 제출을 지시하지는 않는다.
공유 예산, 공통 decay 가격, EfficiencyAdam과 재측정 local residual은
명시적인 변경이다. Native 문장, 손실 reduction, 요청별 종료, ridge,
$C_0$, history와 actual key refresh를 유지한다.

수학 검증은 group-$\ell_1$ 사영의 feasibility와 사영 subproblem 해,
native z 단일 층의 epsilon을 포함한 optimizer 및 clamp parity,
0-norm subgradient, 요청별 종료ㆍfreezeㆍcounter를 확인한다.
Writer 검증은 terminal post-injection 목표, actual key와 current-hidden
refresh, zero-plan 층의 inherited mismatch, native cast, final-key
history append와 rollback을 확인한다. KKT 수치 진단은 convergence
또는 admission gate가 아니다. 이 명세는 검증 통과나 실제 모델의 성능
수치를 선행 주장하지 않는다.

최소 기록 항목은 요청ㆍ층별 계획 norm과 share, 예산 사용률, active set,
$\gamma$ 궤적, 사영 threshold, 구별된 KKT multiplier와 잔차,
candidate 평가ㆍupdate 수와 종료 사유이다. Writer에서는 자기 계획과
inherited mismatch, actual mean-key 및 canonical local 작용,
native cast 이후 realized share를 기록한다. Logical candidate 수,
실제 forward/backward 수, 시간과 메모리는 별도로 보고한다.
계획 후보마다 writer geometry를 재계산하지 않는다. Native prompt는
한 번 tokenize하고, entry anchor와 KL teacher 및 첫 쓰기 층 이전의
고정 prefix는 의미가 동일한 범위에서 재사용할 수 있다. 필요한 loss
위치의 full-vocabulary head, 안전한 오른쪽 padding 제거와 평가
forward 안의 terminal target capture를 사용한다. Contextㆍ어휘ㆍ평가ㆍ
candidate 수를 줄이는 근사로 대체하지 않으며 adapter parity를 확인한다.

편집 평가는 strict ACC, token micro ACC, prompt macro ACC를
rewriteㆍparaphraseㆍneighborhood RS/PS/NS와 함께 보고하며 각 분모와
scoring 정의를 명시한다. Teacher-forced strict ACC를 자유 생성
exact match로 부르지 않는다. 편집 직후 cohort 성능,
후속 batch에 따른 retention과 이웃 true/new NLL도 포함한다.
Matched MEMIT-H 비교는 전체 pipeline의 효과를, 같은 예산과 writer를
사용하는 고정 층 계획 비교는 요청별 배분의 기여를 검사한다.
이 평가값으로 candidate, 층 또는 계수를 고르지 않는다.
배분 다양성 자체가 성능 향상이나 novelty의 증거는 아니며,
locality와 retention의 개선은 관측 결과가 있어야 주장할 수 있다.

Reference, replay, pulse, 추가 old-cost나 warm-up을 도입하지 않는다.
본 문서는 방법 명세이며 실행 승인, dispatch 또는 GPU 실험 결과를
포함하지 않는다. 구현과 실험의 구체적인 profile 및 수치 tolerance는
동반 계약 문서에 기록한다.
\end{document}
